\section{Related Work}
\label{sec:related}

\paragraph{Post-hoc attribution.}
Grad-CAM~\cite{selvaraju2017gradcam}, integrated gradients~\cite{sundararajan2017ig},
SHAP~\cite{lundberg2017shap} and attention rollout~\cite{abnar2020rollout} reconstruct a saliency
map after the prediction. Their reliability has itself been questioned: popular methods
can fail basic sanity checks~\cite{adebayo2018sanity}, and the analogous audit found attention
weights a poor guide to what a model used~\cite{jain2019attention}. We apply that tradition to
intrinsic rather than post-hoc explanation. On dense prediction the extension proceeds
through heuristics, and the
output is an input-\emph{sensitivity} map rather than a concept assignment. We use four such
methods as baselines under an identical protocol, scored by the deletion and insertion curves
introduced with RISE~\cite{petsiuk2018rise}. Their central limitation for our purposes is
structural rather than empirical: a saliency map is a readout, so there is nothing in it to
intervene on. The causal test in Sec.~\ref{sec:protocol} simply cannot be run on them.

\paragraph{Intrinsic interpretability and concept bottlenecks.}
CBMs~\cite{koh2020cbm} and ProtoPNet~\cite{chen2019protopnet} make the explanation part of the
computation, for scalar-output classification. Concepts have also been probed in networks not built around
them~\cite{kim2018tcav,bau2017dissection}; those methods measure alignment between an internal
signal and a concept, precisely the quantity we show is insufficient. Neither addresses dense prediction, where the
explanation must be per-voxel. Critically for this paper, the standard validation in this line is
\emph{concept accuracy or alignment}---does the concept layer predict the annotated
concept---and not whether the downstream head consults it. That is the gap this paper
measures. Reported failure modes for CBMs concern what the concept vector \emph{encodes}---whether
it carries task information beyond the named concepts. Ours is complementary and more basic:
the concept vector may encode the concepts perfectly and still not be read, because the
information bypasses the bottleneck through the \emph{architecture} rather than through
the concept encoding.

\paragraph{Mixture-of-experts routing.}
Sparse MoE~\cite{shazeer2017moe,fedus2022switch,riquelme2021vmoe} introduced routing for
conditional computation. There, routing is causally decisive because the selected expert's output
\emph{is} the layer output. Our results show that transplanting the mechanism into a
skip-connected encoder--decoder silently voids that property---which is exactly the assumption an
interpretability claim would rest on. We inherited the intuition without checking that the
conditions supporting it still held.
